% !TEX spellcheck = en_US



% ~~~~~~~~~~~~~~~~~~~~~~~~~~~~~~~~~~~~~~~ V
% (NC) 2026 Haluk Bingol
% github.com/halukbingol/LaTeX-Templates
%
% Licensees may copy, distribute, display, and perform the work and make derivative works 
% and remixes based on it only for non-commercial purposes. 
% ~~~~~~~~~~~~~~~~~~~~~~~~~~~~~~~~~~~~~~~ A




% =====================================================================
%  Regular Expressions in Java -- Examples: Every Concept as a Real Program
%  ---------------------------------------------------------------------
%  Built on hblib/Template/hblibTemplate.tex with the libraries in ../../hblib/.
%  Programs and their real outputs are in codeoutput/:
%     codeoutput/<Name>.java  a program (classes without main() are only compiled)
%     codeoutput/<Name>.sh    a terminal session, run in a copy of codeoutput/<Name>/
%     codeoutput/<Name>.txt   what it printed, made by:  bash run_examples.sh
% =====================================================================

\documentclass[aspectratio=169,10pt,compress]{beamer}

% ~~~~~~~~~~~~~~~~~~~~~~~~~~~~~~~~~~~~~~~~~~~~~~~~~~~~~~~~~~~~~~~~~ V hblib
	\usepackage{../../hblib/hblibPresentation} % lib presentation: theme, i/N, \hSection, algorithm2e, ...
	\usepackage{../../hblib/hblibCodeoutput} % lib code + output: \lstcode, \lstterminal, ...
	\usepackage{../../hblib/hblibDatastructures} % lib data structures: \hString, \hStack, \hQueue
% ~~~~~~~~~~~~~~~~~~~~~~~~~~~~~~~~~~~~~~~~~~~~~~~~~~~~~~~~~~~~~~~~~ A hblib


% theme, i/N, fonts (T1, lmodern), booktabs, xcolor, array (L{width} columns), listings, tikz,
% algorithm2e (\small captions, \Fn) and the color gitred come from the hblib libraries




% +++++++++++++++++++++++++++++++++++++++ V specific
\DontPrintSemicolon % in algorithms
\newcommand{\win}{\textbf{Windows}}
\newcommand{\mac}{\textbf{macOS}}
% +++++++++++++++++++++++++++++++++++++++ A specific

% ---------------------------------------------------------------------
\title{Regular Expressions in Java}
\subtitle{Examples: Every Concept as a Real Program}
\author[Bingol]{Haluk O. Bingol}

\institute[FCIS]{
    Faculty of Computer and Information Sciences\\
    Yeditepe University
}
\date{
	{\tiny \hbTimeStamp}
}

\begin{document}




% ~~~~~~~~~~~~~~~~~~~~~~~~~~~~~~~~~~~~~~ V frm
\begin{frame}[t,fragile]
  \titlepage
\end{frame}
% ~~~~~~~~~~~~~~~~~~~~~~~~~~~~~~~~~~~~~~ A frm




% ~~~~~~~~~~~~~~~~~~~~~~~~~~~~~~~~~~~~~~ V frm
\begin{frame}[t,fragile] \frametitle{Outline}
  \begin{columns}[T]
    \begin{column}{0.48\textwidth}
      \tableofcontents[sections={1-3}, hideallsubsections]
    \end{column}
    \begin{column}{0.48\textwidth}
      \tableofcontents[sections={4-6}, hideallsubsections]
    \end{column}
  \end{columns}
\end{frame}
% ~~~~~~~~~~~~~~~~~~~~~~~~~~~~~~~~~~~~~~ A frm




% ~~~~~~~~~~~~~~~~~~~~~~~~~~~~~~~~~~~~~~ V frm
\begin{frame}[t,fragile] \frametitle{How to use this deck}
  \begin{itemize}\small
    \item Each concept of the concepts deck, as a small program; the output appears on the second click.
    \item All outputs are real (JDK 21), made by \texttt{bash run\_examples.sh}.
    \item Run one program yourself:
  \end{itemize}
  \begin{columns}[T]
    \begin{column}{0.48\textwidth}
      \begin{block}{Windows (Command Prompt)}


% +++++++++++++++++++++++++++++++++++++++ V lst
%: -Lst.
\begin{lstlisting}[style=codesnippet, language={}, basicstyle=\ttfamily\scriptsize]
cd regex-5-examples\codeoutput
javac Table.java
java Table
\end{lstlisting}
% +++++++++++++++++++++++++++++++++++++++ A lst


      \end{block}
    \end{column}
    \begin{column}{0.48\textwidth}
      \begin{block}{macOS (Terminal)}


% +++++++++++++++++++++++++++++++++++++++ V lst
%: -Lst.
\begin{lstlisting}[style=codesnippet, language=bash, basicstyle=\ttfamily\scriptsize]
cd regex-5-examples/codeoutput
javac Table.java
java Table
\end{lstlisting}
% +++++++++++++++++++++++++++++++++++++++ A lst


      \end{block}
    \end{column}
  \end{columns}

  \small For the Turkish letters on Windows, use Windows Terminal or \texttt{chcp 65001} first.
\end{frame}
% ~~~~~~~~~~~~~~~~~~~~~~~~~~~~~~~~~~~~~~ A frm

% ~~~~~~~~~~~~~~~~~~~~~~~~~~~~~~~~~~~~~~ V
\hSection[Blocks]{Building Blocks}
% ~~~~~~~~~~~~~~~~~~~~~~~~~~~~~~~~~~~~~~ A




% ~~~~~~~~~~~~~~~~~~~~~~~~~~~~~~~~~~~~~~ V frm
\begin{frame}[t,fragile] \frametitle{Which regex matches which text?}
  \lstcoderightstep[0.56][basicstyle=\ttfamily\tiny][basicstyle=\ttfamily\tiny]{codeoutput/Table.java}{Java}{codeoutput/Table.txt}

  \vspace{0.2em}
  \small \texttt{matches}: the \emph{whole} text. \texttt{.} one character, \texttt{*} zero or more, \texttt{+} one or more, \texttt{?} zero or one, \texttt{\{4\}} exactly four, \texttt{|} or.
\end{frame}
% ~~~~~~~~~~~~~~~~~~~~~~~~~~~~~~~~~~~~~~ A frm




% ~~~~~~~~~~~~~~~~~~~~~~~~~~~~~~~~~~~~~~ V frm
\begin{frame}[t,fragile] \frametitle{Anchors, \texttt{find} and \texttt{|}}
  \lstcoderightstep[0.56][basicstyle=\ttfamily\tiny][basicstyle=\ttfamily\tiny]{codeoutput/Anchors.java}{Java}{codeoutput/Anchors.txt}

  \vspace{0.2em}
  \small \texttt{find}: anywhere (F: found). \texttt{\textasciicircum cat|dog\$} is ``starts with cat'' or ``ends with dog''. \texttt{\textbackslash b}: a word boundary. \texttt{matches} needs the whole text.
\end{frame}
% ~~~~~~~~~~~~~~~~~~~~~~~~~~~~~~~~~~~~~~ A frm




% ~~~~~~~~~~~~~~~~~~~~~~~~~~~~~~~~~~~~~~ V frm
\begin{frame}[t,fragile] \frametitle{Groups take a match apart}
  \lstcoderightstep[0.6][basicstyle=\ttfamily\tiny][basicstyle=\ttfamily\tiny]{codeoutput/Dates.java}{Java}{codeoutput/Dates.txt}

  \vspace{0.2em}
  \small \texttt{group(1)}: the first parentheses. In \texttt{replaceAll}, \texttt{\$3-\$2-\$1} puts the parts in a new order.
\end{frame}
% ~~~~~~~~~~~~~~~~~~~~~~~~~~~~~~~~~~~~~~ A frm

% ~~~~~~~~~~~~~~~~~~~~~~~~~~~~~~~~~~~~~~ V
\hSection[Java]{Regexes in Java Source}
% ~~~~~~~~~~~~~~~~~~~~~~~~~~~~~~~~~~~~~~ A




% ~~~~~~~~~~~~~~~~~~~~~~~~~~~~~~~~~~~~~~ V frm
\begin{frame}[t,fragile] \frametitle{Two languages in one line}
  \lstcoderightstep[0.6][basicstyle=\ttfamily\tiny][basicstyle=\ttfamily\tiny]{codeoutput/Escapes.java}{Java}{codeoutput/Escapes.txt}

  \vspace{0.2em}
  \small \texttt{split(".")}: every character is a separator, nothing is left. To match one backslash, the regex is \texttt{\textbackslash\textbackslash}, the Java literal \texttt{"\textbackslash\textbackslash\textbackslash\textbackslash"}.
\end{frame}
% ~~~~~~~~~~~~~~~~~~~~~~~~~~~~~~~~~~~~~~ A frm




% ~~~~~~~~~~~~~~~~~~~~~~~~~~~~~~~~~~~~~~ V frm
\begin{frame}[t,fragile] \frametitle{Compile once, find many: word counts}
  \lstcoderightstep[0.62][basicstyle=\ttfamily\tiny][basicstyle=\ttfamily\tiny]{codeoutput/Words.java}{Java}{codeoutput/Words.txt}

  \vspace{0.2em}
  \small \texttt{\textbackslash p\{L\}+}: Turkish words in one piece; \texttt{\textbackslash w+} breaks them at \texttt{ç}, \texttt{ş}, \texttt{ı}. Lower-case with \texttt{tr}: \texttt{Çay} and \texttt{çay} count together.
\end{frame}
% ~~~~~~~~~~~~~~~~~~~~~~~~~~~~~~~~~~~~~~ A frm

% ~~~~~~~~~~~~~~~~~~~~~~~~~~~~~~~~~~~~~~ V
\hSection[Engines]{How Matching Works}
% ~~~~~~~~~~~~~~~~~~~~~~~~~~~~~~~~~~~~~~ A




% ~~~~~~~~~~~~~~~~~~~~~~~~~~~~~~~~~~~~~~ V frm
\begin{frame}[t,fragile] \frametitle{A finite automaton as a table}
  \lstcoderightstep[0.52][basicstyle=\ttfamily\tiny][basicstyle=\ttfamily\tiny]{codeoutput/Dfa.java}{Java}{codeoutput/Dfa.txt}

  \vspace{0.2em}
  \small The DFA of the concepts deck for \texttt{(a|b)*abb}: one table look-up per character, never back.
\end{frame}
% ~~~~~~~~~~~~~~~~~~~~~~~~~~~~~~~~~~~~~~ A frm




% ~~~~~~~~~~~~~~~~~~~~~~~~~~~~~~~~~~~~~~ V frm
\begin{frame}[t,fragile] \frametitle{Greedy, lazy, exact}
  \lstcoderightstep[0.56][basicstyle=\ttfamily\tiny][basicstyle=\ttfamily\tiny]{codeoutput/Greedy.java}{Java}{codeoutput/Greedy.txt}

  \vspace{0.2em}
  \small \texttt{.*} runs to the \emph{last} \texttt{</b>}; \texttt{.*?} stops at the first; \texttt{[\textasciicircum <]*} cannot pass a \texttt{<}: one way only, nothing to try again.
\end{frame}
% ~~~~~~~~~~~~~~~~~~~~~~~~~~~~~~~~~~~~~~ A frm




% ~~~~~~~~~~~~~~~~~~~~~~~~~~~~~~~~~~~~~~ V frm
\begin{frame}[t,fragile] \frametitle{A backtracking matcher in 35 lines}
  \begin{columns}[T]
    \begin{column}{0.49\textwidth}
      \lstcode[basicstyle=\ttfamily\tiny, lastline=21]{codeoutput/Mini.java}{Java}
    \end{column}
    \begin{column}{0.49\textwidth}
      \lstcode[basicstyle=\ttfamily\tiny, firstline=22, firstnumber=22]{codeoutput/Mini.java}{Java}
    \end{column}
  \end{columns}
\end{frame}
% ~~~~~~~~~~~~~~~~~~~~~~~~~~~~~~~~~~~~~~ A frm




% ~~~~~~~~~~~~~~~~~~~~~~~~~~~~~~~~~~~~~~ V frm
\begin{frame}[t,fragile] \frametitle{Watching it backtrack}
  \begin{columns}[T]
    \begin{column}{0.46\textwidth}
      \lstoutput[basicstyle=\ttfamily\scriptsize]{codeoutput/Mini.txt}
    \end{column}
    \begin{column}{0.50\textwidth}
      \begin{itemize}\small
        \item \texttt{a*} first takes all three \texttt{a}s (greedy): the rest \texttt{ab\$} is tried at 3 and fails.
        \item Give one back: \texttt{ab\$} at 2 fits, then \texttt{b\$}, then \texttt{\$} at the end.
        \item Each \texttt{try} line is one call of \texttt{here}: 5 steps.
        \item \texttt{java.util.regex} works the same way, with many more features.
      \end{itemize}
    \end{column}
  \end{columns}
\end{frame}
% ~~~~~~~~~~~~~~~~~~~~~~~~~~~~~~~~~~~~~~ A frm




% ~~~~~~~~~~~~~~~~~~~~~~~~~~~~~~~~~~~~~~ V frm
\begin{frame}[t,fragile] \frametitle{Counting the steps}
  \begin{columns}[T]
    \begin{column}{0.49\textwidth}
      \lstcode[basicstyle=\ttfamily\tiny, lastline=17]{codeoutput/Steps.java}{Java}
    \end{column}
    \begin{column}{0.49\textwidth}
      \lstcode[basicstyle=\ttfamily\tiny, firstline=18, firstnumber=18]{codeoutput/Steps.java}{Java}
    \end{column}
  \end{columns}
\end{frame}
% ~~~~~~~~~~~~~~~~~~~~~~~~~~~~~~~~~~~~~~ A frm




% ~~~~~~~~~~~~~~~~~~~~~~~~~~~~~~~~~~~~~~ V frm
\begin{frame}[t,fragile] \frametitle{Choices multiply}
  \begin{columns}[T]
    \begin{column}{0.5\textwidth}
      \lstoutput[basicstyle=\ttfamily\scriptsize]{codeoutput/Steps.txt}
    \end{column}
    \begin{column}{0.46\textwidth}
      \begin{itemize}\small
        \item Each line: a failing match on $n$ letters \texttt{a} (no \texttt{c}).
        \item One star: about $n$ steps. Two stars: about $n^{2}/2$. Four stars: about $n^{4}/24$.
        \item The engine tries every way to share the \texttt{a}s among the stars.
        \item Nested stars such as \texttt{(a*)*c} can have $2^{n}$ ways; the senior deck shows when Java pays for that.
      \end{itemize}
    \end{column}
  \end{columns}
\end{frame}
% ~~~~~~~~~~~~~~~~~~~~~~~~~~~~~~~~~~~~~~ A frm

% ~~~~~~~~~~~~~~~~~~~~~~~~~~~~~~~~~~~~~~ V
\hSection[Validate]{Validating Input}
% ~~~~~~~~~~~~~~~~~~~~~~~~~~~~~~~~~~~~~~ A




% ~~~~~~~~~~~~~~~~~~~~~~~~~~~~~~~~~~~~~~ V frm
\begin{frame}[t,fragile] \frametitle{Student numbers and e-mails}
  \lstcoderightstep[0.52][basicstyle=\ttfamily\tiny][basicstyle=\ttfamily\tiny]{codeoutput/Check.java}{Java}{codeoutput/Check.txt}

  \vspace{0.2em}
  \small Input from a file: three lines. \texttt{Omer} fails: \texttt{[a-z]} has no capitals. Compiled once, used for every line.
\end{frame}
% ~~~~~~~~~~~~~~~~~~~~~~~~~~~~~~~~~~~~~~ A frm

% ~~~~~~~~~~~~~~~~~~~~~~~~~~~~~~~~~~~~~~ V
\hSection[Limits]{Form, not Meaning}
% ~~~~~~~~~~~~~~~~~~~~~~~~~~~~~~~~~~~~~~ A




% ~~~~~~~~~~~~~~~~~~~~~~~~~~~~~~~~~~~~~~ V frm
\begin{frame}[t,fragile] \frametitle{A regex checks the form only}
  \lstcoderightstep[0.64][basicstyle=\ttfamily\tiny][basicstyle=\ttfamily\tiny]{codeoutput/Meaning.java}{Java}{codeoutput/Meaning.txt}

  \vspace{0.2em}
  \small \texttt{31.02.2026}: right form, no date. \texttt{LocalDate} (\texttt{STRICT}) knows the calendar.
\end{frame}
% ~~~~~~~~~~~~~~~~~~~~~~~~~~~~~~~~~~~~~~ A frm

% ~~~~~~~~~~~~~~~~~~~~~~~~~~~~~~~~~~~~~~ V
\hSection[Summary]{Summary}
% ~~~~~~~~~~~~~~~~~~~~~~~~~~~~~~~~~~~~~~ A




% ~~~~~~~~~~~~~~~~~~~~~~~~~~~~~~~~~~~~~~ V frm
\begin{frame}[t,fragile] \frametitle{Summary}
  \begin{center}\small
  \begin{tabular}{@{}L{0.3\textwidth}L{0.2\textwidth}L{0.42\textwidth}@{}}
    \toprule
    \textbf{concept}                    & \textbf{program}        & \textbf{lesson} \\
    \midrule
    building blocks                     & \texttt{Table}           & \texttt{matches} = the whole text \\
    anchors, alternation                & \texttt{Anchors}         & \texttt{|} binds weakest; \texttt{find} = anywhere \\
    groups                              & \texttt{Dates}           & \texttt{group(i)}, \texttt{\$i} \\
    two languages                       & \texttt{Escapes}         & every \texttt{\textbackslash} twice in Java \\
    \texttt{Pattern} and \texttt{Matcher} & \texttt{Words}         & compile once; \texttt{\textbackslash p\{L\}} for Turkish \\
    automaton                           & \texttt{Dfa}             & one step per character \\
    backtracking                        & \texttt{Mini}, \texttt{Steps} & greedy, give back; choices multiply \\
    validation                          & \texttt{Check}           & whole-text match of each field \\
    form vs meaning                     & \texttt{Meaning}         & a parser for the value \\
    \bottomrule
  \end{tabular}
  \end{center}
\end{frame}
% ~~~~~~~~~~~~~~~~~~~~~~~~~~~~~~~~~~~~~~ A frm




% ~~~~~~~~~~~~~~~~~~~~~~~~~~~~~~~~~~~~~~ V frm
\begin{frame}[t,fragile] \frametitle{Exercises}
  \begin{enumerate}\small
    \item Add \texttt{[ct]at}, \texttt{c[\textasciicircum a]t} and \texttt{(ca)+t} to \texttt{Table}; predict the rows first.
    \item Change \texttt{Dates} to accept \texttt{2026-11-18} as well and to print all dates as \texttt{18 Kasım 2026}.
    \item Extend \texttt{Mini} with \texttt{+} and \texttt{?}; count its steps on \texttt{a*a*a*c}.
    \item Change \texttt{Dfa} so that it accepts the texts that \emph{contain} \texttt{abb}.
    \item In \texttt{Check}, also read the name and accept Turkish capitals; then reject \texttt{20231101005} if it appears twice.
  \end{enumerate}
\end{frame}
% ~~~~~~~~~~~~~~~~~~~~~~~~~~~~~~~~~~~~~~ A frm




% ~~~~~~~~~~~~~~~~~~~~~~~~~~~~~~~~~~~~~~ V frm
\begin{frame}[t,fragile]
	\begin{center}
		\vspace{4cm}
		\hRemark{\Huge Questions?}
	\end{center}
\end{frame}
% ~~~~~~~~~~~~~~~~~~~~~~~~~~~~~~~~~~~~~~ A frm



\end{document}
